\section{Related Work on Coding Agents}

A range of agent designs have been proposed to improve LLM-based software engineering. SWE-Agent~\citep{sweagent} is first proposed as a foundational system showing that an LLM paired with a small tool set (file editing, command execution, testing, etc.) can iteratively interact with real repositories to resolve issues. Subsequent work refines this paradigm in different directions. Live-SWE-Agent~\citep{xia2025live} studies test-time self-evolution, adapting strategies and occasionally updating prompts, tools, or configuration mid-run based on partial progress. Satori-SWE~\citep{satoriswe} explores population-based evolution at inference time, evolving multiple agent instances or candidates to improve sample efficiency via systematic refinement. In contrast, Agentless~\citep{xia2024agentless} argues for reducing agent complexity by replacing the open-ended loop with a fixed three-stage pipeline (localization, patch generation, test-case generation), achieving strong results on SWE-Bench Lite. Beyond academic prototypes, open-source platforms also shape the ecosystem: OpenHands~\citep{openhands} provides a community toolkit with a unified API for file I/O and code execution and implements a ReAct-style planner over popular base models. More related work can be found at Appendix \ref{sec.more_rw}.
